
\documentclass[letterpaper,11pt]{article}
\oddsidemargin -1.0cm \textwidth 16.5cm

\usepackage[T1]{fontenc}
\usepackage[utf8]{inputenc}
%\usepackage[spanish]{babel}
\usepackage{amsfonts,setspace}
\usepackage{amsmath}
\usepackage{amssymb, amsmath, amsthm}
\usepackage{comment}
\usepackage{amssymb}
\usepackage{dsfont}
\usepackage{anysize}
\usepackage{algorithmicx}
\usepackage{algpseudocode}
\usepackage{multicol}
\usepackage{enumerate}
\usepackage{graphicx}
\usepackage{float}
\usepackage[left=2cm,top=2cm,right=2cm, bottom=2cm]{geometry}
\usepackage{fancyhdr}
\usepackage{algorithm}
\usepackage{algpseudocode}
\pagestyle{fancy}
\usepackage{color}

\theoremstyle{plain}
\newtheorem{teo}{Theorem}
\newtheorem{lem}{Lemma}
\newtheorem{prop}{Prop}
\newtheorem{cor}{Corolary}
\theoremstyle{definition}
\newtheorem{defi}{Definition}
\newtheorem{eje}{Example}
\newtheorem{obs}{Observation}

\newcommand{\R}{\mathbb{R}}
\newcommand{\C}{\mathbb{C}}
\newcommand{\N}{\mathbb{N}}

\usepackage[hidelinks]{hyperref}

\begin{document}

\fancyhead[L]{\itshape{Department of Computer Science}}
\fancyhead[R]{\itshape{University of British Columbia}}

\begin{minipage}{11.5 cm}
\begin{flushleft}
\hspace*{-0.6cm}\textbf{CPSC 320 Intermediate Algorithm Design and Analysis.}\\
\hspace*{-0.6cm}\textbf{Professor: Abner Turkieltaub}
\end{flushleft}\end{minipage}

\begin{center}
	\LARGE \bf{Homework 2}\\
\end{center}

	This assignment is due \textbf{Friday, October 9 at 23:59 PM}. Late submissions will be accepted until Sunday October 11 at 23:59, but penalized: -5 points for Saturday and -15 for Sunday submissions. Please follow these guidelines:
	\begin{itemize}
		\item Prepare your solution using \LaTeX ~and submit a pdf file. Easiest will be to submit using
		the .tex file provided. For questions where you need to select an option, simply change \verb~\square~ ($\square$) to \verb~\blacksquare~ ($\blacksquare$).
		\item Start your solution after our \verb~\color{blue}~ command. Your solution should appear in blue, making it a lot easier for TAs to find what you wrote.
		\item Start each problem on a new page, using the same numbering and ordering as this assignment handout.
		You can use \verb~\newpage~ to start a new page.
		\item Submit the assignment through GradeScope. Your group must make  a \textbf{single} submission via one
		group member's account, marking all other group members in that submission \textbf{using GradeScope's interface}.
		\item  When uploading to Gradescope, link each question with the page(s) of your pdf containing your solution.
	\end{itemize}
	
	Before we begin, a few notes on pseudocode throughout CPSC 320: Your pseudocode should
	communicate your algorithm clearly, concisely, correctly, and without irrelevant detail.
	Reasonable use of plain English is fine in such pseudocode. You should envision your
	audience as a capable CPSC 320 student unfamiliar with the problem you are solving. You may \textbf{neither} include what we consider to be irrelevant coding details \textbf{nor} assume that we understand any particular programming language\footnote{So, for example, do not write \texttt{\#include <iostream>} at the start of your pseudocode, and avoid language-specific notation like C/C++/Java's ternary
	(question-mark-colon) operator.}.\\
	
	Remember also to \textbf{justify/explain  your answers}. We understand that gauging how	much justification to provide can be tricky. Inevitably, judgment is applied by both	student and grader as to how much is enough, or too much, and it's frustrating for all concerned when judgments don't align. Justifications/explanations need not be long or formal, but should be clear and specific (referring back to lines of pseudocode, for example). Proofs should be a bit more formal.\\
	
	On the plus side, if you choose an incorrect answer when selecting an option but your reasoning shows partial understanding, you might get more marks than if no justification is provided. And the effort you put in writing down the justification will hopefully help you gain deeper understanding and may well help you converge to the right selection
	:).
	
	\vspace{.1in}
	
	Ok, time to get started...
	
	\newpage
	
\section{Slow as a logarithm.}
In ``Algorithm Design'' by Kleinberg and Tardos you can find (2.8, page 41) the following useful fact, stating roughly that logarithmic functions are upper-bounded by polynomial functions, more precisely:\\
$$\textbf{Fact:} \text{ for every $b > 1$ and every $x > 0$, we have $\log_b n \in O(n^x)$.}$$
\begin{enumerate}[a)]
\item (10 pts.) Use the fact above to prove the following stronger assertion:
$$
\text{For every $b > 1$ and every $x > 0$, we have $\log_b n \in o(n^x)$.}
$$\\        
\color{blue}
\color{black}
\item (10 pts.) Use part a) to show that $\sqrt{n} \in o(n/\log^{12} n)$. The log is base 10, and $\log^{k} n = (\log n)^k$.
\\
        \color{blue}
        \color{black}
	\end{enumerate}
\newpage
\noindent {\Large\bfseries 1. Slow as a logarithm.}
{\color{blue}
    \begin{enumerate}[a)]
    \item
    To prove that $\log_b n \in o(n^x)$, we need to prove that the $\lim_{n \to \infty} \frac{\log_b n}{n^x} = 0$. And for that, we can leverage the Squeeze Theorem. For the Squeeze Theorem, we need to establish a lower bound and an upper bound. For the lower bound, we will just use 0 (since we want to prove the limit equals 0). For the upper bound, it can be conveniently derived by the fact stated in the question, along with the condition $x > 0$. The fact states that $\log_b n \in O(n^x)$, meaning that there exists some constant $c > 0$, and $n_0$, such that $\log_b n \leq c \cdot n^x$ for all $n \geq n_0$. Since $x > 0$ we can assume that $x \in \mathbb{R}_{>0}$, with $\mathbb{R}_{>0}$ being an uncountable set. This is an important fact because then we know that $\forall x \in \mathbb{R}_{>0}, \exists\alpha \in \mathbb{R}_{>0} : 0 < \alpha < x$ (which is just a fancy way of saying for every $x > 0$, there is some value $\alpha$ that is smaller than $x$ but still greater than $0$). So, we can set up the Squeeze Theorem as such:
    \[\lim_{n \to \infty} 0 \leq \frac{\log_b n}{n^x} \leq \frac{c \cdot n^{\alpha}}{n^x}\]
    \[ = \lim_{n \to \infty} 0 \leq \frac{\log_b n}{n^x} \leq \frac{c}{n^{x - \alpha}}\]
    Since we know that $c$ is constant, even if $x-\alpha$ is tiny, as $n$ goes to $\infty$, the denominator will still get really large. Therefore:
    \[ = \lim_{n \to \infty} 0 \leq \frac{\log_b n}{n^x} \leq 0\]
    By the Squeeze Theorem, since the lower bound and the upper bound both converge to 0, the middle term ($\frac{\log_b n}{n^x}$) must equal 0.
    \[\lim_{n \to \infty} \frac{\log_b n}{n^x} = 0\]
    Hence, as previously stated since $\frac{\log_b n}{n^x} = 0$, we can no denote that for every $b > 1$ and every $x > 0$, we have $\log_b n \in o(n^x)$
    \item
    To show that $\sqrt{n} \in o(n/\log^{12} n)$, we need to show that the $\lim_{n \to \infty} \frac{\sqrt{n}}{\frac{n}{\log^{12} n}} = 0$. First step in showing that would be to clean up the fraction.
    \[\lim_{n \to \infty} \frac{\sqrt{n}}{\frac{n}{\log^{12} n}} = \lim_{n \to \infty} \frac{\log^{12} n}{\sqrt{n}} = \lim_{n \to \infty} \frac{(\log n)^{12}}{\sqrt{n}}\]
    Now, to use part a) in the solution, we want to convert that fraction into the form $\frac{\log_b n}{n^x}$.
    \[\lim_{n \to \infty} \frac{(\log n)^{12}}{\sqrt{n}} = \lim_{n \to \infty} \frac{(\log n)^{12}}{n^{12x}} = \lim_{n \to \infty} \left(\frac{\log n}{n^{x}}\right)^{12}\]
    For the above equalities, we need to make sure that $\sqrt{n} = n^{12x}$ for some $x > 0$. In this case, we want $x = \frac{1}{24}$. We can then rewrite it with $x$ filled in:
    \[\lim_{n \to \infty} \left(\frac{\log n}{n^{\frac{1}{24}}}\right)^{12}\]
    By part a), since the limit above satisfies the form $\frac{\log_b n}{n^x}$ and $\frac{1}{24} > 0$, we know that the term in the base goes towards 0 as $n \to \infty$.
    \[\lim_{n \to \infty} \left(\frac{\log n}{n^{\frac{1}{24}}}\right)^{12} = (\lim_{n \to \infty} \frac{\log n}{n^{\frac{1}{24}}})^{12} = 0^{12} = 0\]
    Therefore we have shown that $\sqrt{n} \in o(n/\log^{12} n)$, through utilizing the proof in a) to specify the stronger assertion.
    \end{enumerate}
}
\newpage		
\section{I wanna be the very best}
In the game ``Pokemon Go'', there is a pvp\footnote{Player vs player.} feature. In this feature, each player selects a team of 3 different pokemons to fight a match against another player. Assume that each pokemon can have one of three roles: ``Lead'', ``Safe swap'' or ``Closer''. We say that a team is \emph{balanced} if it has the three roles\footnote{In case you are curious, you start a match with 2 shields that can be used to block attacks. The lead should be strong at the beginning of the match when both players have shields. The save swap should be versatile and have few hard counters, so when the lead is unfavourable, you can switch to it. The closer is strong at the end of the game when shields are down. All of this is irrelevant to answer the question.}.   
\begin{enumerate}[a)]
    \item (5 pts.) We say that a valid team is just a set of three different pokemons. If you have $n$ pokemons to choose from, how many valid teams are there? \\
    \color{blue}
    %Answer here
    \color{black}
    \item (10 pts.) We want to print a list with all the balanced teams. Consider a brute force approach where you look at every valid team and you print it if it is balanced. Write a pseudo-code to implement this idea using a nested for loop structure, and find a tight $O$-bound for the worst case runtime. Your input is an array of pokemon $[p_1,p_2,\ldots,p_n]$ and you can query the role of any pokemon using $p_i.\text{role()}$ in constant time. Justify your runtime. \\
    \color{blue}
    %Answer here
    \color{black}
    \item (5 pts.) Usually brute force is too slow, but in this case it is not that bad. Assuming that printing one team takes constant time, show that there are instances for which any algorithm that must print all balanced teams has $\Omega(n^3)$ runtime. \\
    \color{blue}
    %Answer here
    \color{black}
    \item (5 pts.) Define the following sets:
    \begin{align*}
        &L= \{p_i : p_i.\text{role()} = \text{Lead}\} \\
        &S= \{p_i : p_i.\text{role()} = \text{Safe Swap}\} \\
        &C= \{p_i : p_i.\text{role()} = \text{Closer}\} 
    \end{align*}
    Write the number of balanced teams as a function of $|L|, \ |S|, \ |R|$. \\
     \color{blue}
    %Answer here
    \color{black}
    \item (10 pts.) In general, you express the runtime as a function of the input size. Sometimes you can analyze the runtime in terms of other parameters to have a more detailed answer. First analyze the runtime of your previous brute force approach in terms of $|L|, \ |S|$ and $|R|$. Then describe a more efficient algorithm in terms of $|L|, \ |S|$ and $|R|$. For your second algorithm, what is the best and worst possible runtime in terms of $n$?  
     \color{blue}
    %Answer here
    \color{black}
\end{enumerate}
\newpage
\noindent {\Large\bfseries 2. I wanna be the very best}
{\color{blue}
\begin{enumerate}[a)]
    \item
    If we say that a valid team is a \textbf{set} of $3$ different pokemons, then the number of valid teams would be $\binom{n}{3}$ or in plain english n choose 3 $(pkm)$. By definition, this is equal to $\frac{(n)_3}{3!} = \frac{n \cdot (n-1) \cdot (n-2)}{3!}$. In this case you would divide by $3!$ to account for double counting (i.e. order does not matter). So, with $\binom{n}{3}$, we would consider teams $(a, b, c)$ and $(a, c, b)$ to be equal. So in total we would have
    \begin{align*}
        \binom{n}{3} &= \frac{n!}{3!\,(n-3)!} = \frac{n!}{6\,(n-3)!} = \frac{n \cdot (n-1)(n-2)}{6} \\[6pt]
        &= \frac{(n^2 - n)(n-2)}{6} = \frac{n^3 - 2n^2 - n^2 + 2n}{6} = \frac{n^3 - 3n^2 + 2n}{6} \text{ total teams}
        \end{align*}
    \item
    The brute force approach is unaware of any pokemon roles, so you would need to first gather all the valid teams (set of size $\binom n3$), then from those pick out the ones that are balanced. In order to find all the valid teams, it would involve a triple-nested loop, where the inner loop moves the fastest, and every nested loop depends on the loop that contains it. For filtering, there are a couple options, but the easiest way to check would just be if all of them have distinct roles. No loop, accumulator, or boolean flags would be required. Since there are 3 slots and 3 roles, we know that if they are all distinct, the team must be balanced. Therefore, we would choose to print it. The pseudocode would look something like this:
    \begin{algorithm}[H]
    \caption{Brute Force Balanced Teams}\label{alg:cap}
    \begin{algorithmic}[1] % [1] enables line numbering for every line
    \Procedure{PrintBalancedTeams}{$P$}
        \State teams = []
        \For{each pokemon i in P}
            \For{each pokemon j = i+1 to n}
                \For{each pokemon k = j+1 to n}
                    \State Append (i, j, k) to teams
                \EndFor
            \EndFor
        \EndFor
        \For{\textbf{each} team $t$ \textbf{in} \textit{teams}}
        \State $r_0, r_1, r_2 \gets t[0].\text{role}(),\, t[1].\text{role}(),\, t[2].\text{role}()$
        \If{$r_0 \neq r_1$ \textbf{and} $r_0 \neq r_2$ \textbf{and} $r_1 \neq r_2$}
            \State Print $t$
        \EndIf
    \EndFor
    \EndProcedure
    \end{algorithmic}
    \end{algorithm}
    To analyze the runtime, we want to look at the triple-nested for loop and the filtering for loop. For the triple-nested for loop, since each of the inner loops depend on the one that contains it (e.g. j = i + 1), the runtime would follow a summation pattern for every set of nested loops. So if we just look at the loops corresponding to pokemon $j$ and pokemon $k$, we would get a runtime of:
    \[\sum_{j=i+1}^n(n-j)\]
    That sum describes the runtime of the inner two loops based on the value of $i$. Then, to account for the outermost loop, you would just sum the runtimes from $i = 1$ to $i = n$.
    \[\sum_{i=1}^n\sum_{j=i+1}^n(n-j)\]
    Now, to figure out the approximate bound of that summation, we're going to need some clever substitutions and the summation formula for $i$ and $i^2$.
    \[\text{Let} ~n-j = m\]
    \[\text{Upper bound:} ~n - (i+1) = n - i - 1\]
    \[\text{Lower bound:} ~n - n = 0\]
    \[\sum_{j=i+1}^n(n-j) = \sum_{m=0}^{n-i-1}m = \frac{(n-i-1)(n-i)}{2}\]
    \[\sum_{i=1}^n\sum_{j=i+1}^n(n-j) = \sum_{i=1}^n\frac{(n-i-1)(n-i)}{2}\]
    \[\text{Let} ~k = n-i\]
    \[\sum_{i=1}^{n}\sum_{j=i+1}^{n}(n-j) = \sum_{k=0}^{n-1}\frac{(k-1)k}{2} = \sum_{k=0}^{n-1}\frac{k^2 - k}{2} = \frac 12\sum_{k=0}^{n-1}(k^2 - k)\]
    \[\frac 12\left(\sum_{k=1}^{n-1}k^2 - \sum_{k=1}^{n-1} k\right) = \frac{1}{2}\left(\frac{(n-1)n(2n-1)}{6} - \frac{(n-1)n}{2}\right)\]
    \[= \frac{n(n-1)(n-2)}{6}\]
    So, with the calculations above, we can conclude that the tight $O$-bound for the triple-nested loop is $O(n^3)$. For the filtering for loop, it is straightforward since we already know how many valid teams exist from part a). That loop would run $\binom n3$ times, which has a tight $O$-bound of $O(n^3)$. Lastly we need to check the roles of all the pokemon in the team, which is 6 queries of constant time, plus printing the balanced team and appending the valid teams, all operations of constant time, and therefore they have a tight $O$-bound of $O(1)$. So with these in mind, the tight $O$-bound of the entire algorithm is $O(n^3)$.
    \item
    An instance where regardless of the algorithm, it would print all the balanced team in $\Omega(n^3)$ runtime would be the case where all role sets have size $\frac n3$. In the case of brute force, it would run in $\Omega(n^3)$ regardless of the shape of the input. This is because you don't know each pokemon's role, meaning you would have to check every pokemon against every other pokemon. In the case of a more efficient algorithm where you already know all pokemons' roles, it still would run in $\Omega(n^3)$ with this particular instance. This is because for each of the $\frac n3$ leads, you would have to check it against the $\frac n3$ safe swaps, and the $\frac n3$ closers, giving $\frac n3 \cdot \frac n3 \cdot \frac n3 = \frac{n^3}{27}$ balanced teams. Since printing a team is a constant operation, printing all of them would run in $\Omega(\frac{n^3}{27}) \in \Omega(n^3)$.
    \item 
    With all the roles segregated into sets, the number of balanced teams would simply be the product of all the set sizes ($|L||S||C|$). In this case, we wouldn't have to divide by $3!$ like we did in part a) because of the fact that $|L|$, $|S|$, and $|C|$ disjoint sets (we can safely assume this from the question "each pokemon can have one of three roles"). Therefore, for all the slots, you wouldn't have to worry about a closer being in the lead spot or something of that nature. All that's to say that we would not have to account for double counting. 
    Hence $Balanced Teams = |L| \cdot |S| \cdot |R|$.
    \item
    Analyzing the runtime of the brute force approach in terms of $|L|$, $|S|$, and $|C|$, we would conclude that the runtime is $O((|L| + |S| + |C|)^3)$. This is because we iterate through every element in the total set of pokemon to get all valid teams, and then check for all balanced teams. This is just equivalent to $O(n^3)$ since we know that all sets are disjoint.
    The efficient way to print out all balanced teams would first start with segregating the pokemons by role. So, we would need to iterate through the list $1$ time. Now that we have all our lead, safe swap, and closer pokemons, we would have $3$ for loops nested inside one another, where for every lead, it would iterate through all safe swaps, which would iterate through all closers. For every closer, we would print the associated balanced team. So, the pseudocode would look something like this:
    \begin{algorithm}[H]
    \caption{Efficient Balanced Teams}\label{alg:cap}
    \begin{algorithmic}[1] % [1] enables line numbering for every line
    \Procedure{PrintBalancedTeams}{$P$}
        \State lead = []
        \State swap = []
        \State closer = []
        \For{each pokemon p in P}
            \If{p.role() = Lead}
                \State Append p to lead
            \ElsIf{p.role() = Safe Swap}
                \State Append p to swap
            \Else
                \State Append p to closer
            \EndIf
        \EndFor
        \For{each pokemon l in lead}
            \For{each pokemon s in swap}
                \For{each pokemon c in closer}
                    \State Print (l,s,c)
                \EndFor
            \EndFor
        \EndFor
    \EndProcedure
    \end{algorithmic}
    \end{algorithm}
    For this second algorithm, the runtime would be $O(|L||S||C| + n)$. The best case runtime would be $O(n)$. This is because the initial for loop to segregate by role would be $O(n)$, then for the triple-nested for loop, an input with something like 1 lead, 1 safe swap, and $n-2$ closers would allow it to run in $O(n)$. We can safely assume this because we know that checking role and printing are constant operations. And only the inner most loop would run $n-2$ times which is $O(n)$. Now for the worst case runtime, it would be $O(n^3)$. This would come as a result of an input like all sets being of size $\frac n3$. The segregation would still be $O(n)$, but the triple-nested for loop would instead be $O(\frac{n^3}{27}) \in O(n^3)$. Since the $n^3$ dominates, the worst case would be $O(n^3)$.
    \end{enumerate}
}
\newpage

\section{Stirling}
The Stirling approximation is a famous and useful formula that approximates the factorial. One way to express it is as follows:
$$\ln (n!) = n\ln n - n+O(\ln n).$$
In other words, $\ln(n!)$ is approximately equal to $n\ln n - n$ plus something bounded by a  logarithmic term. In particular, we can say that $\ln(n!) \in \Theta(n \ln n)$. In this question we will prove this simplified approximation (without using Stirling).  
\begin{enumerate}[a)]
    \item (5 pts.) Prove that $\ln (n!) \in O(n\ln n)$. \\
    \color{blue}

    Choose $c = 1$, $n_0 = 1$
    \begin{align*}
    \ln(n!) &= \ln(n \cdot (n-1)!) \\
    &= \ln(n) + \ln((n-1)!) \\
    &= \ln(n) + \ln(n-1) + \dots + \ln(1) && \text{there are } n \text{ terms} \\
    &\le \ln(n) + \ln(n) + \dots + \ln(n) && \text{since } \ln(n) \text{ is an increasing function} \\
    &= n \cdot \ln(n) \\
    &\le c \cdot n \ln(n)
    \end{align*}
     
    \noindent therefore $\exists c \in \mathbb{R}^+, \exists n_0 \in \mathbb{N}, \forall n \ge n_0, \ln(n!) \le c \cdot n \ln(n)$
     
    \noindent which is the definition of $\ln(n!) \in O(n \ln n)$


    
    \color{black}
    \item (10 pts.) Using induction, prove that $(n!)^2\geq n^n \ \forall n\in \mathbb{N}$ . \\
    \emph{Hint:} it may be useful to prove by induction the following property first: $\forall n \in \mathbb{N}, \ n^n \geq (n+1)^{n-1}$.\\
    \color{blue}

    \noindent proving the hint first: \\
    $\forall n \in \mathbb{N}, n^n \ge (n+1)^{n-1}$
     
    \noindent base case $n = 1$
    \begin{gather*}
    1^1 \ge (1+1)^{1-1} \\
    1 \ge 2^0 \\
    1 \ge 1
    \end{gather*}
     
    \noindent Assume for inductive hypothesis
    \[
    k^k \ge (k+1)^{k-1}
    \]
    for $1 \le k$ and $k \in \mathbb{N}$
     
    \noindent from the inductive hypothesis
    \begin{align*}
    k^k &\ge (k+1)^{k-1} \\
    k^k &\ge \frac{(k+1)^k}{k+1} \\
    \frac{k^k}{(k+1)^k} &\ge \frac{1}{k+1} \\
    \left(\frac{k}{k+1}\right)^k &\ge \frac{1}{k+1} \\
    \left(\frac{k+1}{k+2}\right)^k > \left(\frac{k}{k+1}\right)^k &\ge \frac{1}{k+1} \quad (*) \\
    \left(\frac{k+1}{k+2}\right)^k &\ge \frac{1}{k+1} \\
    \frac{(k+1)^k}{(k+2)^k} &\ge \frac{1}{k+1} \\
    (k+1)^k &\ge \frac{(k+2)^k}{k+1} \\
    (k+1)^k (k+1) &\ge (k+2)^k \\
    (k+1)^{k+1} &\ge (k+2)^k \\
    (k+1)^{k+1} &\ge (k+1+1)^{k+1-1}
    \end{align*}
     
    \noindent to explain $(*)$, if we put $\frac{k+1}{k+2}$ and $\frac{k}{k+1}$ into a ratio, we end up with
    \[
    \frac{k^2 + 2k + 1}{k^2 + 2k}
    \]
    showcasing that $\frac{k+1}{k+2}$ is strictly greater than $\frac{k}{k+1}$
     
    \noindent therefore by mathematical induction, we have
    \[
    \forall n \in \mathbb{N}, n^n \ge (n+1)^{n-1}
    \]
     
    \noindent Next, proving $(n!)^2 \ge n^n, \forall n \in \mathbb{N}$
     
    \noindent base case, $n = 1$
    \[
    (1!)^2 = 1 \ge 1 = 1^1
    \]
     
    \noindent For inductive hypothesis assume that $(k!)^2 \ge k^k, \forall k \in \mathbb{N}, k \ge 1$
    \begin{align*}
    ((k+1)!)^2 &= ((k+1)\, k!)^2 \\
    &= (k+1)^2 (k!)^2 \\
    &\ge (k+1)^2 \cdot k^k && \text{by the inductive hypothesis} \\
    &= (k+1) \cdot (k+1) \cdot k^k \\
    &\ge (k+1) \cdot (k+1) \cdot (k+1)^{k-1} && \text{by using the hint} \\
    &= (k+1) \cdot (k+1)^k \\
    &= (k+1)^{k+1}
    \end{align*}
     
    \noindent therefore by mathematical induction we have
    \[
    (n!)^2 \ge n^n, \forall n \in \mathbb{N}
    \]
     
    \hfill $\blacksquare$
    
    \color{black}
    \item (10 pts.) Use the previous result to conclude that $\ln (n!) \in \Theta(n\ln n)$\\
    \color{blue}

    \noindent from part a, we have $\ln(n!) \in O(n \ln(n))$
     
    \noindent we now seek to prove $\ln(n!) \in \Omega(n \ln(n))$
     
    \noindent choose $c = \frac{1}{2}$ and $n_0 = 1$,
    \begin{align*}
    (n!)^2 &\ge n^n && \text{from part b} \\
    \ln((n!)^2) &\ge \ln(n^n) && \text{since } \ln \text{ is an increasing function and } n \ge n_0 \ge 1 \\
    2\ln(n!) &\ge \ln(n^n) \\
    \ln(n!) &\ge \frac{1}{2} \ln(n^n) && \text{divide both sides by } 2 \\
    &&& \text{then we get} \\
    \ln(n!) &\ge \frac{1}{2} n \ln(n) \\
    &= c\, n \ln(n)
    \end{align*}
     
    \noindent therefore $\exists c \in \mathbb{R}^+, \exists n_0 \in \mathbb{N}, \forall n \ge n_0, c \cdot n \ln(n) \le \ln(n!)$
     
    \noindent which is the definition of $\ln(n!) \in \Omega(n \ln(n))$
     
    \bigskip
     
    \noindent and since we have both $\ln(n!) \in \Omega(n \ln(n))$ and $\ln(n!) \in O(n \ln(n))$
     
    \noindent we can conclude $\ln(n!) \in \Theta(n \ln(n))$
 

    
    \color{black}
    \item (5 pts.) Either prove or provide a counterexample for the following statement: if $h$ is an increasing function and $h(f(n))\in \Theta(h(g(n)))$, then $f(n)\in \Theta(g(n))$. \\
    \color{blue}

    For $h$, we can use $\ln x$ as shown in part a) as a strictly increasing function over its domain. Now, a counterexample would be if you defined $f(n) = n^3$ and $g(n) = n^2$.
    \[h(f(n))\in \Theta(h(g(n)))\]
    \[h(f(n)) \leq c \cdot h(g(n)) ~\text{and} ~h(f(n)) \geq c \cdot h(g(n))\]
    \[h(n^3) \leq c \cdot h(n^2) ~\text{and} ~h(n^3) \geq c \cdot h(n^2)\]
    \[\ln(n^3) \leq c \cdot \ln(n^2) ~\text{and} ~\ln(n^3) \geq c \cdot \ln(n^2)\]
    By the logarithm law shown in part a), the inequalities above can be rewritten as such:
    \[3\ln(n) \leq c \cdot 2\ln(n) ~\text{and} ~3\ln(n) \geq c \cdot 2\ln(n)\]
    Since with $O$ and $\Omega$ we can disregard constants, we can rewrite it as: 
    \[\ln(n) \leq \ln(n) ~\text{and} ~\ln(n) \geq \ln(n)\]
    This means that $h(f(n))\in \Theta(h(g(n)))$. But, for $f(n)\in \Theta(g(n))$ on the other hand, the inequalities would look like this:
    \[n^3 \leq c \cdot n^2 ~\text{and} ~n^3 \geq c \cdot n^2\]
    We can see that for the inequalities above only the right side holds (i.e. $f(n)\in \Omega(g(n))$). That means that $f(n)\notin \Theta(g(n))$.
    
    \color{black}
\end{enumerate}
\newpage

% \noindent {\Large\bfseries 3. Stirling}
% {\color{blue}
    % \begin{enumerate}[a)]
    % \item
    % To prove that $\ln(n!) \in O(n\ln n)$, we need to show that $\ln(n!) \leq c \cdot n\ln n$ for some constant $c > 0$. First, through the use of logarithm laws we know that $\ln(a^n) = n\ln(a)$. So, we can rewrite $n\ln n$ as $\ln(n^n)$. So, we would then be left with this inequality to prove:
    % \[\ln(n!) \leq \ln(n^n)\]
    % Now, for strictly increasing functions, we can use the property of strict monotonicity and exponentiate both sides with base $e$ to get the same inequality but with only the terms within the parentheses (i.e. $n!$ and $n^n$). Now, we can show that $\ln x$ is strictly increasing over its domain by taking the first derivative and checking if it is always positive ($> 0$).
    % \[\frac{dy}{dx}\ln x = \frac 1x\]
    % \[\forall x > 0, ~\frac 1x > 0\]
    % \[\therefore n! \leq n^n \]
    % Now, the inequality above is much nicer to prove. And for this question we are assuming that $n \geq 0$ since factorials are only defined on non-negative integers. For this proof we are going to expand both terms and compare factors.
    % \[n \cdot (n-1) \cdot (n-2) \cdot ... \cdot 1 \leq n \cdot n \cdot n \cdot ... \cdot n\]
    % \[(n - 1) < n, ~(n-2) < n, ~... ,~1 < n\]
    % Since we established the domain being only non-negative integers, we can safely say that comparing individual factors like that will hold. Therefore $\ln(n!) \leq \ln(n^n) = \ln(n!) \leq c \cdot n\ln n$, meaning $\ln(n!) \in O(n\ln n)$.
    % \item
    % \item
    % \item
    % For $h$, we can use $\ln x$ as shown in part a) as a strictly increasing function over its domain. Now, a counterexample would be if you defined $f(n) = n^3$ and $g(n) = n^2$.
    % \[h(f(n))\in \Theta(h(g(n)))\]
    % \[h(f(n)) \leq c \cdot h(g(n)) ~\text{and} ~h(f(n)) \geq c \cdot h(g(n))\]
    % \[h(n^3) \leq c \cdot h(n^2) ~\text{and} ~h(n^3) \geq c \cdot h(n^2)\]
    % \[\ln(n^3) \leq c \cdot \ln(n^2) ~\text{and} ~\ln(n^3) \geq c \cdot \ln(n^2)\]
    % By the logarithm law shown in part a), the inequalities above can be rewritten as such:
    % \[3\ln(n) \leq c \cdot 2\ln(n) ~\text{and} ~3\ln(n) \geq c \cdot 2\ln(n)\]
    % Since with $O$ and $\Omega$ we can disregard constants, we can rewrite it as: 
    % \[\ln(n) \leq \ln(n) ~\text{and} ~\ln(n) \geq \ln(n)\]
    % This means that $h(f(n))\in \Theta(h(g(n)))$. But, for $f(n)\in \Theta(g(n))$ on the other hand, the inequalities would look like this:
    % \[n^3 \leq c \cdot n^2 ~\text{and} ~n^3 \geq c \cdot n^2\]
    % We can see that for the inequalities above only the right side holds (i.e. $f(n)\in \Omega(g(n))$). That means that $f(n)\notin \Theta(g(n))$.
    % \end{enumerate}
% }

\newpage
\section{It's not the destination, it's the journey.}
In planet Namek, all the alien kids have now their school for the upcoming school year. For parents, the question now is what is the fastest way to get to school. For parents with only one kid this is trivial since they can use noodle maps to find the shortest path from their home to the school\footnote{In noodle map you can enter two locations $A$ and $B$ and it tells you the shortest path from $A$ to $B$.}. Consider a parent of twins in different schools, this parent wants to find the shortest route to drop both kids and then return home\footnote{In Namek, all parent can work from home.}. 
\begin{enumerate}[a)]
\item (5 pts.) For each of the following scenarios, explain how they can find the shortest route using noodle maps the least number of times possible.
\begin{itemize}
    \item Case 1: distances are symmetric (distance from $A$ to $B$ is the same as the distance from $B$ to $A$).
    \item Case 2: distances are asymmetric\footnote{Maybe there's a school in the top of a mountain and you have to take some long stairs to get there, but then when you leave you can go down a slide.}. \\

\end{itemize}
\item (5 pts.) Many algorithms for shortest path (like Dijkstra) actually find the shortest path from one vertex to all the others. Imagine you are a developer for noodle maps and you want to implement a query where the user (a parent) inputs their home and the two schools and it shows them the shortest route starting from their home, passing by the two schools and returning home. Assume that as a developer, when you run noodle maps to find a shortest path from $A$ to $B$, you can actually access the shortest paths from $A$ to every other location. Answer the questions from the previous part under this new hypothesis. \\

\item (5 pts.) Answer the questions from the previous two parts, but now under the assumption that work from home is not allowed, so parents must go to their work after dropping both kids.
\end{enumerate}

\newpage
\noindent {\Large\bfseries 4. It's not the destination, it's the journey.}
{\color{blue}


\begin{quote}
    Let: \\
    \texttt{H} be home \\
    \texttt{A} be school A \\
    \texttt{B} be school B \\
    \texttt{W} be work
    \end{quote}

\begin{enumerate}[a)]

    \item
        For Case 1, they can find the shortest route using noodle maps 3 times due to symmetry. As stated by the problem, the distance from A to B is the same as from B to A. So whether the parent checked $\text{dist}(H, A)$, $\text{dist}(A, B)$, and $\text{dist}(B, H)$, it would be symmetric to the case where the parent checked $\text{dist}(H, B)$, $\text{dist}(B, A)$, and $\text{dist}(A, H)$. Therefore, we would just have to check one of those instances to get the shortest route. 

        For Case 2, they can find the shortest route using noodle maps 6 times due to asymmetry. Since the distance from A to B is not the same as the distance from B to A, in this case, to find the shortest route, the parent has to check on the noodle maps which route is faster: driving to school A first or to school B first. These are denoted as either $\text{dist}(H, A)$, $\text{dist}(A, B)$, and $\text{dist}(B, H)$, or $\text{dist}(H, B)$, $\text{dist}(B, A)$, and $\text{dist}(A, H)$, which, as previously stated, is a total of 6 times that noodle maps is utilized.

    \item  
        Given that as a developer, when we run noodle maps to find a shortest path from A to B, we get access to all the shortest paths from A to every location, for Case 1, we would only need use noodle maps 2 times. This is because running from (H, A) gives implicitly (H, B), and we just need to run it again from (A,B), to get all the distances to be able to drop of both children and arrive home. 
        
        For Case 2, holding the additonal edges revealed as a developer it we would need to use noodle maps 3 times. This is because Case 2 is assymetric so we require  we need all 6 edges to calculate the shortest path:
        
        \begin{verbatim}
        (H, A)
        (A, H)
        (A, B)
        (B, A)
        (H, B)
        (B, H)
        \end{verbatim}
        
        As previously stated, running noodle maps on \texttt{(H, A)} gives also \texttt{(H, B)}. Then, on \texttt{(A, B)} gives \texttt{(A, H)}, and on \texttt{(B, H)} gives \texttt{(B, A)}. Thus, after 3 uses of noodle maps we would get all 6 edges to be able to calculate the shortest path given each of their distances.
    \item
    
    
    In resolving part a with the notion of part c, that parents cannot work from home, we need 5 uses of noodle maps if symmetric (Case 1). This is because we are no longer returning home, and therefore the final leg is different for each route: we go to one school, then the second school, and then to work, instead of coming back home from the second school. In part a, this is what allowed us to cut the number of searches, since the distances (H, A) and (H, B) were always covered. Hence we need to check the following edges:
    
    \begin{verbatim}
    (H, A)
    (H, B)
    (A, B)
    (A, W)
    (B, W)
    \end{verbatim}

    Checking the 5 edges except (B, A), due to symmetry, allows us to ensure that we choose the cumulative shortest distance in going from home to a school, and from the second school to work, as traveling between schools is always the same due to symmetry.

    On the other hand, for Case 2, we need to check noodle maps 6 times, due to the asymmetry,  one per each edge:
    \begin{verbatim}
    (H, A)
    (H, B)
    (A, B)
    (B, A)
    (A, W)
    (B, W)
    \end{verbatim}

    In the Case 1, we were able to avoid checking (B, A), since the distance from A to B is the same from B to A, but when the paths are no longer symmetric, we need to ensure we have the values of all 6 different possible edges to ensure, we are to correctly identify the edges that cummulative result in the shortest path.
    
    Now working under part notions in b:
    
    We state that in Case 1, we only need to use noodle maps 2 times due to symmetric. In this case similarly to the how we solved part a, we need only five edges, yet we can get them through asking from the edges (A, H) and (B, W), since it gives us all 6 edges: 
    
    \begin{verbatim}
    (A, H) -> (A, B), (A, W), (H, A) - by symmetry
    (B, W) -> (B, A), (B, W), (B, H) = (H, B) by symmetry
    \end{verbatim}

    Therefore we have all the distance from the home to both schools, and to work, in order to be able to find the shortest path.
    
    In the asymmetric case, which is Case 2, we need to use noodle maps 3 times, since we can no longer get all the edges requeried in simply 2 times. This is because previously (A, H), would allow us to get (H, A) due to symmetry, but no we need to query directly to H, since symmetry no longer holds. We would then simply need to query the following:
    
    \texttt{(A, W)} gives \texttt{(A, B)}, 
    
    \texttt{(B, A)} gives, \texttt{(B, W)}. 
    
    \texttt{(H, A)} gives \texttt{(H, B)}

    Then we would have all 6 required edges,

    \begin{verbatim}
        (H, A)
        (H, B)
        (A, W)
        (A, B)
        (B, W)
        (B, A)
    \end{verbatim}

    to calculate what the shortest distance is between home and each school, and then from each school to work.
    
    \color{black}
    \end{enumerate}
    }
\newpage


	\newpage
	    
	\section*{Statement on Collaboration and Use of Resources}
	
	To develop good practices in doing homeworks, citing resources and acknowledging input from others, please complete the following.\\

    None of your answers here will be used against you, so please be honest. Not answering this question will result in a deduction of 2 points.
	
	\begin{enumerate}
        \item Please list the CWLs of all group members. \\
        \color{blue}
        Andrew Li - andrwhli \\
        Sean Yan - dhsyan \\
        Alonso Bastidas Garcia - abastida
        
        \color{black}
		\item List here any books or online sources that you consulted and that led you to the solution of one or more problems (indicate the problem(s) after each resource): \\
		\color{blue}
        \url{https://en.wikipedia.org/wiki/List_of_logarithmic_identities} - Problem 1 \\
        \url{https://en.wikipedia.org/wiki/Squeeze_theorem} - Problem 1 \\
        \url{https://en.wikipedia.org/wiki/Summation} - Problem 2 \\
        
        \color{black}
		\item One or more of us consulted with course staff during office hours:\\
		$\square$ Yes \hspace{.5in} $\blacksquare$ No
		
		\item One or more of us collaborated with other CPSC 320 students outside our group:\\		
		$\square$ Yes \hspace{.5in} $\blacksquare$ No\\
		If yes, please list their name(s) here:
        \color{blue}

        
        \color{black}
		\item One or more of us used an AI to solve one or more questions from this assignment:\\
        $\blacksquare$ Yes \hspace{.5in} $\square$ No\\

        \color{blue}
		If yes, please say which problem(s) was(were) solved using AI and which AI model you used:
        We used Claude Opus 5.5 in some questions. Mainly to only answer 
        us yes or no questions, to tell us if we are on the right path. 
        Any other more explicit usage, was statement in line in the question.

        We also used AI, to format latex. 
        
        \color{black}
        \item One or more of us collaborated with or consulted others outside of CPSC 320:\\
        $\square$ Yes \hspace{.5in} $\blacksquare$ No\\
		If yes, please list their name(s) here:
        \color{blue}

        
        
    \vspace{.1in}

	\end{enumerate}
 	   

\end{document}

